\section{Related Work}



\paragraph{Tool-using LLMs and execution feedback.}
ReAct interleaves language reasoning with environment actions
\citep{yao2023react}, while Toolformer and ToolLLM study how language models
can acquire and invoke external tools \citep{schick2023toolformer,qin2023toollm}.
This line of work primarily addresses action selection and tool-use
capability. PERSIST-ACE starts from an already operational tool-using LLM and
focuses on a later execution problem: a syntactically plausible action can
continue to produce no new evidence, no task-relevant state change, or the
same unresolved functional effect.

\paragraph{Execution recovery and loop handling.}
Early-exit methods reduce repetitive or ineffective behavior by deciding when
to halt, but termination alone does not repair a recoverable task state
\citep{lu2025runaway}. Other approaches generate recovery plans at test time,
as in Reasoning Inception \citep{kim2026rein}, or update the LLM policy so
that it learns from execution errors, as in Fission-GRPO
\citep{zhang2026fission}. PERSIST-ACE takes a narrower and more conservative
route: it does not terminate every loop, invoke an open-ended planner, or
update model parameters at inference time. Instead, it applies at most one
frozen, grounded, deterministic correction when the trigger and safety gates
pass, and preserves the Raw action when grounding is ambiguous or a retry may
still be productive. Unlike early-exit, generative replanning, or policy
optimization, its central object is a selective programmatic escape from a
feedback-visible Functional Loop.
